\documentclass[leqno]{article}
\usepackage{verbatim}
\usepackage{array}
\usepackage{listings}
\usepackage{fancyvrb}
\usepackage{enumitem}

\usepackage[utf8]{inputenc}
\usepackage[T1]{fontenc}
\usepackage{textcomp}
\usepackage{multicol} \usepackage{mathtools}
\usepackage{amsmath}
\usepackage{wrapfig}
\usepackage{amssymb}
\usepackage{amsmath,amsfonts,amssymb,amsthm,epsfig,epstopdf,titling,url,array}
\usepackage{hyperref}
\usepackage{eso-pic}
\usepackage{pgf}
\usepackage{tikz}
\usepackage{tikz-cd}
\usepackage{graphicx}

% figure support
\usepackage{import}
\usepackage{xifthen}
\pdfminorversion=7
\usepackage{pdfpages}
\usepackage{transparent}
\usepackage{xcolor}

% geometry
\usepackage{geometry}
\geometry{a4paper, margin=1in}

% paragraph length
\setlength{\parindent}{0em}
\setlength{\parskip}{1em}

\newtheorem*{theorem}{Theorem}
\newtheorem*{lemma}{Lemma}
\newtheorem*{proposition}{Proposition}
\newtheorem*{definition}{Definition}
\newtheorem*{observation}{Observation}

\DeclareMathOperator{\Hom}{Hom}
\DeclareMathOperator{\Poset}{Poset}
\DeclareMathOperator{\im}{Im}
\DeclareMathOperator{\coker}{Coker}
\DeclareMathOperator{\coim}{Coim}

\newcommand{\com}[1]{\textcolor{red}{#1}}
\newcommand{\incfig}[1]{%
\center
\def\svgwidth{0.9\columnwidth}
\import{./figures/}{#1.pdf_tex}
}
\newcommand{\incimg}[1]{%
\center
\includegraphics[width=0.9\columnwidth]{images/#1}
}
\pdfsuppresswarningpagegroup=1

\title{Entregable Topología Algebraica}
\author{Abel Doñate Muñoz}
\date{}

\begin{document}
\maketitle
\tableofcontents
\newpage

\section{Homología en la categoría de Poset}

Para ver que $C_{\bullet}(P)$ es, en efecto, un complejo, debemos probar que la aplicación $d$ es la aplicación borde, es decir  $d^2 = 0$.
 \begin{align*}
 &d\circ d (x_0\le \cdots \le x_n) = \sum_{i=0}^n (-1)^{i} d(x_0\le  \cdots \le  \hat{x_i} \le \cdots\le x_n) = \\
 &=\sum_{i=0}^n \left( \sum_{j=0}^{i-1} (-1)^{i+j}(x_0\le \cdots\le \hat{x_j}\le \cdots \le \hat{x_i}\le \cdots \le x_n)  - \sum_{j=i+1}^{n} (-1)^{i+j} (x_0\le \cdots\le \hat{x_i}\le \cdots \le \hat{x_j}\le \cdots \le x_n) \right) 
\end{align*}
Pero esta expresión es cero, ya que si fijamos dos indices diferentes $k, l$, vemos que pasa una vez por  $(i, j)=(k, l)$ y otra por  $(j, i) = (k, l)$ con signo cambiado anulándose. Ya que toda cadena es combinación lineal de cadenas con coeficientes en $\mathbb{Z}$, al extender la cadena por linealidad del morfismo $d$ tenemos que  $d^2$ sigue dando  $0$ para toda cadena. Esto prueba que 
% https://q.uiver.app/#q=WzAsNSxbMSwwLCJDX3tuKzF9Il0sWzIsMCwiQ19uIl0sWzMsMCwiQ197bi0xfSJdLFswLDAsIlxcY2RvdHMiXSxbNCwwLCJcXGNkb3RzIl0sWzAsMSwiZCJdLFsxLDIsImQiXSxbMywwLCJkIl0sWzIsNCwiZCJdXQ==
\[\begin{tikzcd}
	\cdots & {C_{n+1}(P)} & {C_n(P)} & {C_{n-1}(P)} & \cdots
	\arrow["d_{n+1}", from=1-2, to=1-3]
	\arrow["d_{n}", from=1-3, to=1-4]
	\arrow["d_{n+2}", from=1-1, to=1-2]
	\arrow["d_{n-1}", from=1-4, to=1-5]
\end{tikzcd}\]
es en efecto un complejo. Al ser un complejo en una categoría abeliana, podemos definir la homología como
\[
H_n(P) = \ker(d_n) / \im (d_{n+1}) := Z_n / B_{n+1}
\]
que esta bien definida gracias a que es un complejo y la imagen queda dentro del núcleo.

c) Se puede comprobar que la homología con $ \le $ y la realizada con $<$ da lugar a grupos de homología equivalentes. Para los ejemplos utilizaremos las cadenas de $<$, ya que son más fáciles de computar.  

Para los ejemplos empezaremos por calcular $H_0$. Observamos que $d_1(x)=0\ \forall x$, por lo que dado el conjunto del poset con cardinal $n$ tenemos que todo el grupo abeliano $\sum a_n x_n$ esta en el nucleo de  $d_1$, por lo que  $Z_0 \cong  \mathbb{Z}^n$. De la misma forma si consideramos la imagen $B_1$, esta esta formada por todas las cadenas de la forma $x-y$ si  $x\le y$, es decir, si están en la misma componente conexa del grafo de Hasse (no importa la dirección, pues si está al reves haremos $y-x$). Por tanto al pasar por el cociente estamos identificando todas las componentes conexas y $H_0 \cong \mathbb{Z}^{\# \{\text{comp. con.}\}}$.

\textbf{Primer poset} 

Por el argumento anterior hemos visto que $H_0 = \mathbb{Z}$. Calculamos ahora $H_1$. El nucleo de una cadena arbitraria de $C_1$ será
\[
d_1(ax_0x_1 + bx_1x_3 + cx_0x_2 + dx_2x_3) = x_1 (a-b) + x_2(c-d) + x_0(-a-c-e) + x_3(b+d+e) = 0
\]
lo que quiere decir
\[
a=b, \quad c=d, a+c+e=0, \quad \implies \ker d_1 = \mathbb{Z}^2
\]

Por otra parte si calculamos la imagen de $C_2$ tenemos
\[
d_2(ax_0x_1x_3 + bx_0x_2x_3) = a(x_1x_3 - x_0x_3 + x_0x_1) + b(x_2x_3-x_0x_3+x_0x_2) \implies \im(d_2) = \mathbb{Z}^2
\]
por tanto $H_1 = 0$. 

Calcular $H_2$ es más sencillo, pues el núcleo ya da cero, ya que si agrupamos los términos de la imagen de un elemento arbitrario como el anterior tenemos 
\[
ax_1x_3 + ax_0x_1 + bx_2x_3 + bx_0x_2 + (-a-b)x_0x_3 = 0 \implies a=b=0 \implies \ker (d_2) = 0
\]
por lo que $H_2 = 0$. 

En definitiva $H_0 = \mathbb{Z}, H_i = 0 \ \forall  i \ge  1$.


\textbf{Segundo poset} 

Como el grafo de Hasse es conexo, entonces $H_0 = \mathbb{Z}$. Calculamos ahora el nucleo de la aplicación $d_1$
\[
d_1(ax_1x_3 + bx_1x_2+cx_0x_2) = ax_3 -cx_0 +(-a-b)x_1 + (b+c)x_2 =0 \implies a=b=c = 0 \implies \ker (d_1)=0
\]
lo que significa que $H_1 = 0$. 

En definitiva $H_0 = \mathbb{Z}, H_i = 0 \ \forall  i \ge  1$.

\textbf{Tercer poset} 

Como el grafo de Hasse es conexo, entonces $H_0 = \mathbb{Z}$. Calculamos ahora el nucleo de la aplicación $d_1$
\begin{align*}
d_1(ax_0x_3 + bx_0x_4 + cx_1x_3+dx_1x_4 + ex_2x_3 + fx_2x_4) = &(-a-b) x_0 + (-c-d)x_1+(-e-f)x_2 +\\
&+ (a+c+e) x_3 + (-b-d-f)x_4 = 0
\end{align*}
Nos salen las ecuaciones
\[
b = -a, \quad d = -c, \quad f = -e, a+c+e = 0 \implies (a,c,e): a+c+e = 0
\]
lo que es claramente isomorfo a $\mathbb{Z}^2$. por otro lado la imagen de $d_2$ es $0$, ya que las cadenas de longitud 3 $C_2 = 0$. Deducimos, por tanto, que $H_1 = \mathbb{Z}^2$.

En definitiva $H_0 = \mathbb{Z}, H_1 = \mathbb{Z}^2, H_i = 0 \ \forall  i \ge  2$ .



\section{Complejos simpliciales abstractos}

a) Para probar que $\Delta (K)$ es complejo de cadenas debemos ver que $\partial_{p}\circ \partial _{p-1} = 0$. Esto se realiza de manera totalmente equivalente al ejercicio 1. Vamos a ver que es un functor
% https://q.uiver.app/#q=WzAsNixbMCwwLCJLIl0sWzEsMCwiTCJdLFsyLDAsIlMiXSxbMCwxLCJcXERlbHRhIEsiXSxbMSwxLCJcXERlbHRhIEwiXSxbMiwxLCJcXERlbHRhIFMiXSxbMCwzLCJGIl0sWzAsMSwiZiJdLFsxLDIsImciXSxbNCw1LCJmXyoiLDJdLFszLDQsImZfKiIsMl0sWzEsNCwiRiJdLFsyLDUsIkYiXV0=
\[\begin{tikzcd}
	K & L & S \\
	{\Delta K} & {\Delta L} & {\Delta S}
	\arrow["f", from=1-1, to=1-2]
	\arrow["F", from=1-1, to=2-1]
	\arrow["g", from=1-2, to=1-3]
	\arrow["F", from=1-2, to=2-2]
	\arrow["F", from=1-3, to=2-3]
	\arrow["{f_*}"', from=2-1, to=2-2]
	\arrow["{f_*}"', from=2-2, to=2-3]
\end{tikzcd}\]
Si la función $f$ envía vértices a vertices, podemos definir la $f_*:S(K) \to S(L)$ que envie elemento a elemento dentro del conjunto. De esta forma el diagrama conmuta. Ver que $F(g \circ f) = g_* \circ  f_*$ se puede hacer partiendo desde un elemento de $\Delta (K)$ y ver que, en efecto, la composición por abajo (la que envía elementos a elementos dentro del conjunto) es exactamente la misma que la composición por arriba aplicando el functor.  

b) Tenemos los siguientes diagramas

\begin{minipage}{0.5\textwidth}
% https://q.uiver.app/#q=WzAsNixbMCwwLCJcXERlbHRhIEsiXSxbMSwwLCJcXERlbHRhIEwiXSxbMiwwLCJcXERlbHRhIFMiXSxbMCwxLCJIX1xcYnVsbGV0KFxcRGVsdGEoSykpIl0sWzEsMSwiSF9cXGJ1bGxldChcXERlbHRhKEwpKSJdLFsyLDEsIkhfXFxidWxsZXQoXFxEZWx0YShTKSkiXSxbMCwzLCJGIl0sWzAsMSwiZiJdLFsxLDIsImciXSxbNCw1LCJnXyoiLDJdLFszLDQsImZfKiIsMl0sWzEsNCwiRiJdLFsyLDUsIkYiXV0=
\[\begin{tikzcd}
	{\Delta (K)} & {\Delta (L)} & {\Delta (S)} \\
	{H_\bullet(\Delta(K))} & {H_\bullet(\Delta(L))} & {H_\bullet(\Delta(S))}
	\arrow["f", from=1-1, to=1-2]
	\arrow["F", from=1-1, to=2-1]
	\arrow["g", from=1-2, to=1-3]
	\arrow["F", from=1-2, to=2-2]
	\arrow["F", from=1-3, to=2-3]
	\arrow["{f_*}"', from=2-1, to=2-2]
	\arrow["{g_*}"', from=2-2, to=2-3]
\end{tikzcd}\]
\end{minipage}
\begin{minipage}{0.5\textwidth}
% https://q.uiver.app/#q=WzAsNixbMCwwLCJcXERlbHRhIEsiXSxbMSwwLCJcXERlbHRhIEwiXSxbMiwwLCJcXERlbHRhIFMiXSxbMCwxLCJIX1xcYnVsbGV0KFxcRGVsdGEoSykpIl0sWzEsMSwiSF9cXGJ1bGxldChcXERlbHRhKEwpKSJdLFsyLDEsIkhfXFxidWxsZXQoXFxEZWx0YShTKSkiXSxbMCwzLCJGIl0sWzAsMSwiZiJdLFsxLDIsImciXSxbNCw1LCJnXyoiLDJdLFszLDQsImZfKiIsMl0sWzEsNCwiRiJdLFsyLDUsIkYiXV0=
\[\begin{tikzcd}
	{\Delta' (K)} & {\Delta' (L)} & {\Delta' (S)} \\
	{H_\bullet(\Delta'(K))} & {H_\bullet(\Delta'(L))} & {H_\bullet(\Delta'(S))}
	\arrow["f", from=1-1, to=1-2]
	\arrow["F", from=1-1, to=2-1]
	\arrow["g", from=1-2, to=1-3]
	\arrow["F", from=1-2, to=2-2]
	\arrow["F", from=1-3, to=2-3]
	\arrow["{f_*}"', from=2-1, to=2-2]
	\arrow["{g_*}"', from=2-2, to=2-3]
\end{tikzcd}\]
\end{minipage}
donde el functor de homología se define como el grupo abeliano resultante $H_n(\Delta (K)) = \ker \partial_n / \im (\partial _{n+1})$. Podemos ver que es functorial, ya que si tenemos una aplicación $f: \Delta(K) \to \Delta(L)$, entonces tenemos la aplicación induida en sus nucleos. Tomando la clase de equivalencia con las imagenes tenemos una aplicación inducida $f_*$ que hace conmutar el diagrama.

Para probar la functorialidad de $\Delta '(K)$ tenemos que hacer un paso extra, que consiste en ver que la aplicación $\partial '_p$ está bien definida, y no importa la permutación que elijamos. Podemos ver esto de manera que si tomamos dos permutaciones arbitrarias del conjunto $(v_0, \dots, v_p)$, entonces la imagen por $\partial _p$ será la misma, pero posiblemente con el signo cambiado, lo que indica que las imágenes son las mismas tras hacer la identificación propuesta.

c, d) Las homologías orientadas serán exactamente igual que las homologías de posets del problema 1. Podemos verlo haciendo la identificación 
 \[[(v_0 , v_1 , \dots , v_n )] \leftrightarrow v_0 < v_1 < \dots < v_n\]
 Esta idenitficación manda la clase de un conjunto de vértices al conjunto ordenado de vértices que es elemento del poset.


Calculamos ahora las homologías orientadas. Por conexión todos los grupos de homología de orden 0 serán $H_0 = \mathbb{Z}$.

\textbf{Primer poset}
Si hacemos la imagen de un elemento genérico de $C_2$ que es $ x = a_1(1, 2, 4) + a_2(1, 4, 2) + a_3(2, 4, 1) + a_4(2, 1, 4) + a_5(4, 1, 2) + a_6(4, 2, 1) + b_1(1, 3, 4) + b_2(1, 4, 3) + b_3(3, 4, 1) + b_4(3, 1, 4) + b_5(4, 1, 3) + b_6(4, 3, 1) $ tenemos 
\begin{align*}
d(x) = &(2,4)(a_1+a_3-a_4) + (4, 2)(a_2+a_6-a_5) + (1,2) (a_1-a_2+a_5) + (2,1) (-a_3+a_4+a_6) +\\
&+(3,4)(b_1+b_3-b_4) + (4, 3)(b_2+b_6-b_5) + (1,3) (b_1-b_2+b_5) + (3,1) (-b_3+b_4+b_6) \\
& + (1,4) (-a_1+a_2+a_4-b_1+b_2+b_4) +(4,1)(a_3+a_5-a_6+b_3+b_5-b_6)
\end{align*}
La imagen de este morfismo es el rango de la matriz $10\times 12$ que define las ecuaciones. Calculando el rango con ayuda del siguiente script de python tenemos que $ \im(d_2) \cong \mathbb{Z}^7$. Por teorema de isomorphismo $\mathbb{Z}^{12}/ \ker (d_2) \cong \im (d_2) \implies \ker (d_2) \cong \mathbb{Z}^5$, por lo que $H_2 = \mathbb{Z}^5$.

\begin{verbatim}
	import numpy as np

	A = np.array([[1,0,1,-1,0,0,   0,0,0,0,0,0],
				  [0,1,0,0,-1,1,   0,0,0,0,0,0],
				  [1,-1,0,0,1,0,   0,0,0,0,0,0],
				  [0,0,-1,1,0,1,   0,0,0,0,0,0],
				  [0,0,0,0,0,0,    1,0,1,-1,0,0],
				  [0,0,0,0,0,0,    0,1,0,0,-1,1],
				  [0,0,0,0,0,0,    1,-1,0,0,1,0],
				  [0,0,0,0,0,0,    0,0,-1,1,0,1],
				  [-1,1,0,1,0,0,   -1,1,0,1,0,0],
				  [0,0,1,0,1,-1,   0,0,1,0,1,-1]])
	
	r = np.linalg.matrix_rank(A)
	print(r)
\end{verbatim}

De manera análoga tenemos el siguiente elemento genérico de $C_1$ que es $x = a_1(1,2) + a_2(2,1) + b_1(1,3)+b_2(3,1) + c_1(1,4) + c_2(4,1) + d_1(2,4)+d_2(4,2)+ e_1(3,4) + e_2(4,3)$. Si hacemos la diferencial vemos que nos sale exactamente lo mismo que en el caso del ejercicio (1), pero con las variables $a=a_1-a_2, b=b_1-b_2$, etc. Por tanto el kernel será el mismo pero con 5 variables libres más. $ \ker (d_1) = \mathbb{Z}^2\times \mathbb{Z}^5 = \mathbb{Z}^7$. Deducimos, entonces que $H_1 = 0$.

En resumen: $H_0 = \mathbb{Z}, H_1 = 0, H_2 = \mathbb{Z}^5, H_{i} = 0 \ \forall  i \ge 3$ 

\textbf{Segundo poset}
Por lo visto en el primer poset podemos utilizar que $ker(d_1) = \mathbb{Z}^0 \times \mathbb{Z}^3 = \mathbb{Z}^3$, donde $\mathbb{Z}^0$ es el cálculo del problema (1) y $\mathbb{Z}^3$ corresponde a las variables libres. Por tanto $H_1 = \mathbb{Z}^3$.

En resumen: $H_0 = \mathbb{Z}, H_1 = \mathbb{Z}^3, H_i = 0 \ \forall  i \le 2$.

\textbf{Tercer poset}
Por lo visto en el primer poset podemos utilizar que $ker(d_1) = \mathbb{Z}^2 \times \mathbb{Z}^6 = \mathbb{Z}^8$, donde $\mathbb{Z}^2$ es el cálculo del problema (1) y $\mathbb{Z}^6$ corresponde a las variables libres. Por tanto $H_1 = \mathbb{Z}^8$.

En resumen: $H_0 = \mathbb{Z}, H_1 = \mathbb{Z}^8, H_i = 0 \ \forall  i \le 2$.











\section{Cilindro}

a) Para probar que es un complejo de cadenas debemos probar que las imagenes del morfismo caen donde toca, lo que es trivial, y que $d_Z\circ d_Z=0$.
\begin{align*}
  & d_Z\circ d_Z (a', a'', b) = d_Z (d_A(a')+a'', -d_A(a''), -f(a'')+d_B(b)) =\\
  =& (d_A(d_A(a')+a'')-d_A(A'')), -d_A(d_A(a'')) , f(d_A(a'')) + d_B(-f(a'')+d_B(b)) ) =\\
  =& (0, 0, f(d_A(a''))-d_B(f(a'')))= (0,0,0)
\end{align*}

\begin{minipage}{0.7\textwidth}
Donde hemos utilizado $d_B^2=d_A^2=0$ y que el diagrama de la derecha conmuta
\end{minipage}
\begin{minipage}{0.3\textwidth}
% https://q.uiver.app/#q=WzAsNCxbMCwwLCJBX24iXSxbMSwwLCJBX3tuLTF9Il0sWzAsMSwiQl9uIl0sWzEsMSwiQl97bi0xfSJdLFswLDIsImYiLDJdLFsxLDMsImYiXSxbMCwxLCJkX0EiXSxbMiwzLCJkX0IiLDJdXQ==
\[\begin{tikzcd}
	{A_n} & {A_{n-1}} \\
	{B_n} & {B_{n-1}}
	\arrow["f"', from=1-1, to=2-1]
	\arrow["f", from=1-2, to=2-2]
	\arrow["{d_A}", from=1-1, to=1-2]
	\arrow["{d_B}"', from=2-1, to=2-2]
\end{tikzcd}\]
\end{minipage}

b) Para probar que $\xi$ es morfismo de cadenas tenemos que probar que el siguiente diagrama conmuta.
% https://q.uiver.app/#q=WzAsOCxbMCwwLCJCX24iXSxbMSwwLCJCX3tuLTF9Il0sWzAsMSwiWl9uIl0sWzEsMSwiWl97bi0xfSJdLFszLDAsImIiXSxbNCwwLCJkX0IoYikiXSxbMywxLCIoMCwwLGIpIl0sWzQsMSwiKDAsIDAsIGRfQihiKSkiXSxbMCwyLCJcXHhpX24iLDJdLFsyLDMsImRfWiIsMl0sWzEsMywiXFx4aV97bi0xfSJdLFswLDEsImRfQiJdLFs0LDUsImRfQiIsMCx7InN0eWxlIjp7InRhaWwiOnsibmFtZSI6Im1hcHMgdG8ifX19XSxbNCw2LCJcXHhpX24iLDIseyJzdHlsZSI6eyJ0YWlsIjp7Im5hbWUiOiJtYXBzIHRvIn19fV0sWzYsNywiZF9aIiwyLHsic3R5bGUiOnsidGFpbCI6eyJuYW1lIjoibWFwcyB0byJ9fX1dLFs1LDcsIlxceGlfe24tMX0iLDAseyJzdHlsZSI6eyJ0YWlsIjp7Im5hbWUiOiJtYXBzIHRvIn19fV1d
\[\begin{tikzcd}
	{B_n} & {B_{n-1}} && b & {d_B(b)} \\
	{Z_n} & {Z_{n-1}} && {(0,0,b)} & {(0, 0, d_B(b))}
	\arrow["{\xi_n}"', from=1-1, to=2-1]
	\arrow["{d_Z}"', from=2-1, to=2-2]
	\arrow["{\xi_{n-1}}", from=1-2, to=2-2]
	\arrow["{d_B}", from=1-1, to=1-2]
	\arrow["{d_B}", maps to, from=1-4, to=1-5]
	\arrow["{\xi_n}"', maps to, from=1-4, to=2-4]
	\arrow["{d_Z}"', maps to, from=2-4, to=2-5]
	\arrow["{\xi_{n-1}}", maps to, from=1-5, to=2-5]
\end{tikzcd}\]
Vemos que, en efecto, conmuta, ya que $d_Z\circ \xi_n = \xi_{n-1}\circ d_B$.

Para probar que $\eta$ es morfismo de cadenas tenemos que probar que el siguiente diagrama conmuta.

% https://q.uiver.app/#q=WzAsOCxbMCwwLCJaX24iXSxbMSwwLCJaX3tuLTF9Il0sWzAsMSwiQl9uIl0sWzEsMSwiQl97bi0xfSJdLFszLDAsIihhJywgYScnLCBiKSJdLFs0LDAsIihkX0EoYScpK2EnJywgLWRfQShhJycpLCAtZihhJycpK2RfQihiKSkiXSxbMywxLCJmKGEnKStiIl0sWzQsMSwiZF9CKGIpICsgZihkX0EoYScpKSA9IGRfQihiKSsgZF9CKGYoYScpKSJdLFswLDIsIlxcZXRhX24iLDJdLFsyLDMsImRfQiIsMl0sWzEsMywiXFxldGFfe24tMX0iXSxbMCwxLCJkX1oiXSxbNCw1LCJkX1oiLDAseyJzdHlsZSI6eyJ0YWlsIjp7Im5hbWUiOiJtYXBzIHRvIn19fV0sWzQsNiwiXFxldGFfbiIsMix7InN0eWxlIjp7InRhaWwiOnsibmFtZSI6Im1hcHMgdG8ifX19XSxbNiw3LCJkX0IiLDIseyJzdHlsZSI6eyJ0YWlsIjp7Im5hbWUiOiJtYXBzIHRvIn19fV0sWzUsNywiXFxldGFfe24tMX0iLDAseyJzdHlsZSI6eyJ0YWlsIjp7Im5hbWUiOiJtYXBzIHRvIn19fV1d
\[\begin{tikzcd}
	{Z_n} & {Z_{n-1}} && {(a', a'', b)} & {(d_A(a')+a'', -d_A(a''), -f(a'')+d_B(b))} \\
	{B_n} & {B_{n-1}} && {f(a')+b} & {d_B(b) + f(d_A(a')) = d_B(b)+ d_B(f(a'))}
	\arrow["{\eta_n}"', from=1-1, to=2-1]
	\arrow["{d_B}"', from=2-1, to=2-2]
	\arrow["{\eta_{n-1}}", from=1-2, to=2-2]
	\arrow["{d_Z}", from=1-1, to=1-2]
	\arrow["{d_Z}", maps to, from=1-4, to=1-5]
	\arrow["{\eta_n}"', maps to, from=1-4, to=2-4]
	\arrow["{d_B}"', maps to, from=2-4, to=2-5]
	\arrow["{\eta_{n-1}}", maps to, from=1-5, to=2-5]
\end{tikzcd}\]
Vemos que, en efecto, conmuta, ya que $d_B\circ \eta_n = \eta_{n-1}\circ d_Z$.

c) Tenemos que ver que $\varphi = \xi \circ \eta \simeq Id_{Z(f)}$ y que $\eta \circ \xi \simeq Id_{B}$. La segunda es evidente, ya que la composición ya da la identidad. Tenemos que probar, pues, que existe una homotopía de cadenas $s_i:Z_i \to Z_{i+1} $ tal que $\varphi - Id_Z = s_{n-1}\circ d_n + d_{n+1}\circ s_n$, donde $\varphi = \xi \circ  \eta $ tal que $\varphi ((a', a'', b))= \eta (f(a')+b) = (0, 0, f(a')+b)$. 

% https://q.uiver.app/#q=WzAsMTAsWzEsMCwiWl97bisxfSJdLFsxLDEsIlpfe24rMX0iXSxbMywwLCJaX3tuLTF9Il0sWzMsMSwiWl97bi0xfSJdLFsyLDAsIloiXSxbMiwxLCJaIl0sWzQsMCwiXFxjZG90cyJdLFs0LDEsIlxcY2RvdHMiXSxbMCwwLCJcXGNkb3RzIl0sWzAsMSwiXFxjZG90cyJdLFs0LDUsIlxcdmFycGhpIl0sWzQsMSwic19uIiwyXSxbMSw1LCJkX3tuKzF9IiwyXSxbMiw1LCJzX3tuLTF9IiwyXSxbNCwyLCJkX24iXSxbMCw0LCJkX3tuKzF9Il0sWzUsMywiZF9uIiwyXSxbMiw2XSxbMyw3XSxbOCwwXSxbOSwxXV0=
\[\begin{tikzcd}
	\cdots & {Z_{n+1}} & Z & {Z_{n-1}} & \cdots \\
	\cdots & {Z_{n+1}} & Z & {Z_{n-1}} & \cdots
	\arrow["\varphi", from=1-3, to=2-3]
	\arrow["{s_n}"', from=1-3, to=2-2]
	\arrow["{d_{n+1}}"', from=2-2, to=2-3]
	\arrow["{s_{n-1}}", from=1-4, to=2-3]
	\arrow["{d_n}", from=1-3, to=1-4]
	\arrow["{d_{n+1}}", from=1-2, to=1-3]
	\arrow["{d_n}"', from=2-3, to=2-4]
	\arrow[from=1-4, to=1-5]
	\arrow[from=2-4, to=2-5]
	\arrow[from=1-1, to=1-2]
	\arrow[from=2-1, to=2-2]
\end{tikzcd}\]

Ahora bien, la forma más natural de definir la aplicación $s:Z_n = A_n\oplus A_{n-1} \oplus B_n\to Z_{n+1}=A_{n+1}\oplus A_{n}\oplus B_{n+1}$ es enviar $(a', a'', b)  \to  (0, \pm a' , 0)$. Vemos que con el signo negativo tenemos la homotopía deseada:
% https://q.uiver.app/#q=WzAsOSxbMSwwLCIoYScsIGEnJywgYikiXSxbMCwxLCIoMCwgLWEnLCAwKSJdLFsyLDEsIihkX0EoYScpK2EnJywgLWRfQShhJycpLCAtZihhJycpK2RfQihiKSkiXSxbMCwyLCIoLWEnLGRfQShhJyksIGYoYScpKSJdLFsyLDIsIigwLCAtZF9BKGEnKS1hJycsIDApIl0sWzEsMiwiKyJdLFszLDAsIigtYScsIC1hJycsIGYoYSkpIl0sWzMsMiwiKC1hJywgLWEnJywgZihhKSkiXSxbMywxLCI9Il0sWzAsMSwic19uIiwyXSxbMCwyLCJkX24iXSxbMiw0LCJzX3tuLTF9IiwyXSxbMSwzLCJkX3tuKzF9Il0sWzAsNiwiXFx2YXJwaGkgLSBJZCJdLFs0LDcsIisiXV0=
\[\begin{tikzcd}
	& {(a', a'', b)} && {(-a', -a'', f(a))} \\
	{(0, -a', 0)} && {(d_A(a')+a'', -d_A(a''), -f(a'')+d_B(b))} & {=} \\
	{(-a',d_A(a'), f(a'))} & {+} & {(0, -d_A(a')-a'', 0)} & {(-a', -a'', f(a))}
	\arrow["{s_n}"', from=1-2, to=2-1]
	\arrow["{d_n}", from=1-2, to=2-3]
	\arrow["{s_{n-1}}"', from=2-3, to=3-3]
	\arrow["{d_{n+1}}", from=2-1, to=3-1]
	\arrow["{\varphi - Id}", from=1-2, to=1-4]
	\arrow["{+}", from=3-3, to=3-4]
\end{tikzcd}\]


\section{Toro cociente}

(a) \com{terminar}

Consideramos la inclusión $(T_f-B,A-B) \hookrightarrow (T_f, A)$. Claramente $\overline{B} \subseteq A^\circ$, por lo que podemos emplear escisión y tenemos $H_{\bullet}(T_f-B, A-B) \cong H_{\bullet}(T_f, A)$. Vemos ahora las siguientes propiedades:

\begin{enumerate}[topsep=-6pt, itemsep=0pt]
	\item $T_f-B \simeq X \times (0,1) \simeq X \times I$
	\item $A-B \simeq X \times \partial I \simeq X \sqcup X$
	\item $A \simeq X_0$ 
\end{enumerate}

La primera es consecuencia de que estamos quitando la parte que identificamos y existe un retracto de deformación de $X \times (\varepsilon , 1-\varepsilon ) \to X \times (0, 1)$. La segunda es consecuencia de que nos quedamos con las partes $X\cdot (0,\varepsilon ) \cup X \cdot (0,\varepsilon )$, y existe un retracto de deformación hacia $X\times  \{0\} \cup X \times \{1\}$. La tercera es simplemente que existe el retracto de deformación $A \to X_0$ que envía $\overline{(x,a)} \to \overline{(x,0)}$. 

Tenemos, por tanto
\begin{align*}
\begin{cases}
H _{\bullet} (T_f-B, A-B) \cong H _{\bullet} (X\times I, X \times \partial I) \\
H _{\bullet}(T_f, A) \cong H _{\bullet} (T_f, X_0)
\end{cases} \implies H _{\bullet} (X \times I, X \times \partial I) \cong H _{\bullet}  (T_f, X_0)
\end{align*}

(b) 
% https://q.uiver.app/#q=WzAsNixbMSwwLCJIX3tuKzF9KFhcXHRpbWVzIEkpIl0sWzIsMCwiSF97bisxfShYXFx0aW1lcyBJLCBYXFx0aW1lcyBcXHBhcnRpYWwgSSkiXSxbMywwLCJIX3tufShYXFx0aW1lc1xccGFydGlhbCBJKSJdLFs0LDAsIkhfbihYXFx0aW1lcyBJKSJdLFswLDAsIlxcY2RvdHMiXSxbNSwwLCJcXGNkb3RzIl0sWzAsMSwial8qID0gMCJdLFsxLDIsIlxccGFydGlhbF8qIiwwLHsic3R5bGUiOnsidGFpbCI6eyJuYW1lIjoiaG9vayIsInNpZGUiOiJ0b3AifX19XSxbMiwzLCJcXGlvdGFfKiIsMCx7InN0eWxlIjp7ImhlYWQiOnsibmFtZSI6ImVwaSJ9fX1dLFs0LDAsIlxcaW90YV8qIl0sWzMsNSwial8qPTAiXV0=
\[\begin{tikzcd}
	\cdots & {H_{n+1}(X\times I)} & {H_{n+1}(X\times I, X\times \partial I)} & {H_{n}(X\times\partial I)} & {H_n(X\times I)} & \cdots
	\arrow["{j_* = 0}", from=1-2, to=1-3]
	\arrow["{\partial_*}", hook, from=1-3, to=1-4]
	\arrow["{\iota_*}", two heads, from=1-4, to=1-5]
	\arrow["{\iota_*}", from=1-1, to=1-2]
	\arrow["{j_*=0}", from=1-5, to=1-6]
\end{tikzcd}\]

Vemos que $\iota _{\bullet}$ es exhaustivo, ya que $H_{\bullet}(X\times \partial I) \cong H _{\bullet}(X \sqcup X) \cong H _{\bullet}(X)\times H _{\bullet}(X)$ y $H _{\bullet}(X\times I) \cong H _{\bullet}(X)$ . Por tanto tenemos que $\ker (j_{\bullet}) = \im (\iota _{\bullet}) = H _{\bullet} (X) \implies j _{\bullet}=0$.

Como $j_* = 0$, tenemos que  $\partial_*$ es monomorfismo y $\iota_*$ es epimorfismo. Por tanto
\[
H_{n+1}(X\times I, X\times \partial I) \cong  \im (\partial_*) \cong  \ker (\iota_*) 
\] 
Pero como $H_n(X\times \partial I) \simeq H_n(X)\times H_n(X)$, ya que está formado por dos copias de los espacios, tenemos que, al ser $\iota_*$ epimorfismo y $X\times I \simeq X$:
\begin{align*}
\iota_*: H_n(X)\times H_n(X) \to H_n(X \times I) \cong  H_n(X) \quad  \Rightarrow \quad \ker (\iota_*) \cong H_n(X) \quad \Rightarrow \quad H_{n+1}(X\times I, X\times \partial I) \cong H_n(X)
\end{align*}
como queríamos ver.

(c)  \com{falta	}

Consideramos la secuancia exacta larga de abajo del diagrama. Sabemos por un lado que 
\[
H _{\bullet} (T_f, X_0) \cong H _{\bullet}(X \times I, X \times \partial I) = \im (\partial _{\bullet})= \ker (\iota _{\bullet}) \cong H _{\bullet} (X)
\]



\end{document}
